%% GENERATED by tools/from_docs.py from stepss-docs. Do not edit by hand:
%% edit the page named below in stepss-docs and regenerate.
%% source: stepss-docs/src/content/docs/models/synchronous-machine-param-conversion.md

\chapter{Synchronous machine parameter conversion}\label{doc:model-sync-conv}

A \texttt{SYNC\_MACH} record can be entered using one of two equivalent parameter formats, selected by the \texttt{TYPE\_MOD} keyword:

\begin{itemize}
\tightlist
\item
  \textbf{\texttt{RL}}: the Park-model inductances and resistances are supplied directly.
\item
  \textbf{\texttt{XT}}: characteristic reactances and open-circuit time constants are supplied; STEPSS/RAMSES converts them internally.
\end{itemize}

Both formats describe the same machine. \texttt{XT} is convenient when data comes from manufacturer datasheets or Kundur-style standard parameters. This page documents exactly how that conversion is done, so it can be reproduced by hand or cross-checked against external simulators (e.g.~Typhoon HIL EMT).

\begin{docsnote}{Common pitfalls}

Two details that break hand calculations against STEPSS:
1. Open-circuit time constants are normalised to \textbf{radians} by \(t_b = 1/(2\pi f_{nom})\) before use in the resistance formulas.
2. The field circuit base is tied to \texttt{IBRATIO} via \(p_{uf} = R_f / \text{IBRATIO}\); a wrong assumption here corrupts \(M_d\), \(L_{\ell f}\), and \(R_f\) after per-unit \ensuremath{\rightarrow} Henry conversion.

\end{docsnote}

\begin{center}\rule{0.5\linewidth}{0.5pt}\end{center}

\section{Authoritative Field Order}\label{model-sync-conv:authoritative-field-order}

From the RAMSES source (\texttt{sync.f90}, \texttt{get\_sync\_mach}):

\begin{Verbatim}
SYNC_MACH NAME BUS FP FQ P Q SNOM PNOM H D IBRATIO
  RL  LL MDU  LLF  LLD1 MQU LLQ1 LLQ2 M N RA RF   RD1  RQ1  RQ2
  XT  LL XD   XPD  XSD  XQ  XPQ  XSQ  M N RA TPD0 TSD0 TPQ0 TSQ0
  EXC <model> ... TOR <model> ... ;
\end{Verbatim}

A rotor circuit the machine does not have is skipped with \texttt{*} in \textbf{both} its reactance/inductance and its resistance/time-constant field.

\begin{center}\rule{0.5\linewidth}{0.5pt}\end{center}

\section{Conversion Algorithm (XT \ensuremath{\rightarrow} RL)}\label{model-sync-conv:conversion-algorithm-xt-rl}

All parameters are in per unit; time constants are entered in seconds and normalised internally.

\textbf{Time base}

\[
t_b = \frac{1}{2\pi f_{nom}}
\]

All time constants (TPD0, TSD0, TPQ0, TSQ0) are divided by \(t_b\) before entering the resistance formulas.

\textbf{Unsaturated mutuals}

\[
M_d^u = X_d - L_\ell, \qquad M_q^u = X_q - L_\ell
\]

\textbf{d axis, two rotor circuits (field + damper, XSD/TSD0 present)}

\[
T'_d = T'_{d0}\frac{X'_d}{X_d}, \qquad T''_d = \frac{X''_d \, T'_{d0} \, T''_{d0}}{X_d \, T'_d}
\]

\[
a = \frac{X_d T'_d T''_d - L_\ell T'_{d0} T''_{d0}}{X_d - L_\ell}, \quad
b = \frac{X_d(T'_d + T''_d) - L_\ell(T'_{d0} + T''_{d0})}{X_d - L_\ell}
\]

\[
c = \frac{X_d(T'_{d0} T''_{d0} - T'_d T''_d)}{(X_d - L_\ell)^2}, \quad
d = \frac{X_d(T'_{d0} + T''_{d0} - T'_d - T''_d)}{(X_d - L_\ell)^2}
\]

\[
T_f = \frac{b + \sqrt{b^2 - 4a}}{2}, \qquad T_{d1} = b - T_f
\]

\[
R_{d1} = \frac{T_f - T_{d1}}{c - d \, T_{d1}}, \qquad R_f = \frac{-R_{d1}}{1 - d \, R_{d1}}
\]

\[
L_{\ell f} = T_f R_f, \qquad L_{\ell d1} = T_{d1} R_{d1}
\]

\textbf{d axis, single rotor circuit (XSD/TSD0 skipped)}

\[
L_{ff} = \frac{(M_d^u)^2}{X_d - X'_d}, \qquad L_{\ell f} = L_{ff} - M_d^u, \qquad R_f = \frac{L_{ff}}{T'_{d0}}
\]

\textbf{q axis} is fully symmetric to the d axis. Two circuits (XPQ/TPQ0 and XSQ/TSQ0) use the same quadratic with \(X_q, L_\ell, T'_{q0}, T''_{q0}, X'_q, X''_q\). Single-circuit fallbacks:

\begin{itemize}
\tightlist
\item
  Transient only (XPQ/TPQ0): \(L_{\ell q1} = M_q^u{}^2/(X_q - X'_q) - M_q^u\), \(R_{q1} = (M_q^u{}^2/(X_q - X'_q))/T'_{q0}\)
\item
  Subtransient only (XSQ/TSQ0): same with \(X''_q, T''_{q0}\), result assigned to \(L_{\ell q2}, R_{q2}\)
\end{itemize}

\textbf{Field base scaling}

\[
p_{uf} = R_f \,/\, \text{IBRATIO}
\]

The field current is then reconstructed internally as \(i_f = (\psi_f - \psi_{ad}) \cdot (R_f/p_{uf}) / L_{\ell f}\).

\begin{docsnote}{Sanity check}

If any of \(R_f, R_{d1}, L_{\ell f}, L_{\ell d1}\) (or q-axis equivalents) turns out negative, the supplied reactances/time constants are physically inconsistent. STEPSS emits an ``unrealistic Park inductances or resistances'' warning.

\end{docsnote}

\begin{center}\rule{0.5\linewidth}{0.5pt}\end{center}

\section{Reference Python Implementation}\label{model-sync-conv:reference-python-implementation}

A standalone Python port of the \texttt{XT} branch is available at
\href{https://github.com/SPS-L/Sync_mach_Octave}{\texttt{Sync\_mach\_Octave}} (Octave) and
as \texttt{ramses\_xt\_to\_park.py} (attached below). It reproduces the algorithm above
including the \(t_b\) normalisation, the quadratic solve, the symmetric q axis,
and \texttt{puf\ =\ RF/IBRATIO}. Pass \texttt{None} for any rotor circuit the machine does not
have.

\begin{Verbatim}
from ramses_xt_to_park import ramses_xt_to_park

p = ramses_xt_to_park(
    fnom=50.0, ibratio=1.0,
    ll=0.15,  ra=0.003,
    xd=1.81,  xpd=0.30, xsd=0.23, tpd0=8.0, tsd0=0.03,
    xq=1.76,  xpq=0.65, xsq=0.25, tpq0=1.0, tsq0=0.07)
# -> llf=0.169902, lld1=0.166338, rf=0.000741, rd1=0.033390, ...
\end{Verbatim}

\begin{center}\rule{0.5\linewidth}{0.5pt}\end{center}

\section{Worked Examples}\label{model-sync-conv:worked-examples}

Inputs (both cases, 50 Hz, IBRATIO = 1): \(L_\ell = 0.15\), \(R_a = 0.003\), \(X_d = 1.81\), \(X'_d = 0.30\), \(X''_d = 0.23\), \(T'_{d0} = 8.0\) s, \(T''_{d0} = 0.03\) s, \(X_q = 1.76\), \(X'_q = 0.65\), \(T'_{q0} = 1.0\) s. The round-rotor case adds \(X''_q = 0.25\), \(T''_{q0} = 0.07\) s.

\begin{small}\begin{longtable}[]{@{}lrr@{}}
\toprule\noalign{}
Parameter & Round rotor (d2/q2) & Single q-damper (d2/q1) \\
\midrule\noalign{}
\endhead
\bottomrule\noalign{}
\endlastfoot
\(M_d^u\) & 1.660000 & 1.660000 \\
\(L_{\ell f}\) & 0.169902 & 0.169902 \\
\(L_{\ell d1}\) & 0.166338 & 0.166338 \\
\(M_q^u\) & 1.610000 & 1.610000 \\
\(L_{\ell q1}\) & 0.928153 & 0.725225 \\
\(L_{\ell q2}\) & 0.120461 & n/a \\
\(R_f\) & 0.000741 & 0.000741 \\
\(R_{d1}\) & 0.033390 & 0.033390 \\
\(R_{q1}\) & 0.009236 & 0.007433 \\
\(R_{q2}\) & 0.028210 & n/a \\
\(p_{uf}\) & 0.000741 & 0.000741 \\
\end{longtable}\end{small}

\begin{center}\rule{0.5\linewidth}{0.5pt}\end{center}

\section{Comparing with an EMT Simulator}\label{model-sync-conv:comparing-with-an-emt-simulator}

When cross-checking STEPSS (RMS/phasor) against an EMT tool such as Typhoon HIL:

\begin{enumerate}
\def\labelenumi{\arabic{enumi}.}
\tightlist
\item
  \textbf{Per unit vs physical units.} STEPSS stays entirely in per unit. Converting to Henry/Ohm for the EMT tool requires a consistent per-unit base on the rotor side; use the same \texttt{IBRATIO} assumption on both sides.
\item
  \textbf{Radians vs seconds.} The \(t_b\) normalisation is internal to STEPSS. Your hand calculation must divide every time constant by \(t_b = 1/(2\pi f_{nom})\) before computing resistances.
\item
  \textbf{Expected post-fault difference.} STEPSS uses the phasor approximation and neglects stator transformer voltages (\(d\psi_d/dt\), \(d\psi_q/dt\)), so it omits the DC-offset and high-frequency current components immediately after a short circuit. An EMT model retains them. Compare the \textbf{slow post-fault envelopes} first: envelope agreement with first-cycle differences indicates a modelling assumption, not a parameter error.
\end{enumerate}

\begin{center}\rule{0.5\linewidth}{0.5pt}\end{center}

\section{See Also}\label{model-sync-conv:see-also}

\begin{itemize}
\tightlist
\item
  \hyperref[doc:model-sync]{Synchronous Machine Model}, equations, per unit system, and \texttt{SYNC\_MACH} record reference
\item
  \href{https://github.com/SPS-L/Sync_mach_Octave}{Octave reference implementation}
\item
  \href{https://thierryvancutsem.github.io/home/elec0047/phasor_approx.pdf}{Phasor approximation}
\item
  \href{https://thierryvancutsem.github.io/home/elec0047/dyn_of_sync_mac.pdf}{Synchronous machine dynamics}
\end{itemize}
